\section{主要模組設計}

\begin{table}[htbp]
    \centering
    \caption{系統模組功能說明}
    \label{tab:module_description}
    \renewcommand{\arraystretch}{1.3}
    \begin{tabularx}{\textwidth}{
        >{\centering\arraybackslash}p{4cm}
        >{\raggedright\arraybackslash}X
    }
        \toprule
        模組名稱 & 功能說明 \\
        \midrule
        \texttt{top} & 系統最上層，連接所有子模組。 \\
        \texttt{audio\_codec\_ctrl} & 設定 WM8731，例如取樣率、輸入輸出模式。 \\
        \texttt{i2c\_config} & 透過 I\textsuperscript{2}C 設定 WM8731 暫存器。 \\
        \texttt{i2s\_rx} & 接收 Audio CODEC 的音訊資料。 \\
        \texttt{i2s\_tx} & 傳送音訊資料至 Audio CODEC 播放。 \\
        \texttt{sram\_controller} & 控制 SRAM 讀寫。 \\
        \texttt{flash\_controller} & 控制 FLASH 讀寫。 \\
        \texttt{record\_controller} & 錄音流程控制。 \\
        \texttt{playback\_controller} & 播放流程控制。 \\
        \texttt{fft\_core} & 對音訊資料進行 FFT 分析。 \\
        \texttt{spectrum\_mapper} & 將 FFT 結果轉成 8 階 LED 高度。 \\
        \texttt{led\_matrix\_driver} & 掃描控制 8$\times$8 LED 矩陣。 \\
        \texttt{lcd\_controller} & 控制 16$\times$2 LCD 顯示。 \\
        \texttt{button\_debounce} & BUTTON 防彈跳。 \\
        \texttt{system\_fsm} & 控制整體模式切換。 \\
        \bottomrule
    \end{tabularx}
\end{table}

\subsection{FFT 與 LED 矩陣顯示設計}

\subsubsection{頻率範圍}

本系統取樣頻率為 22.4 kHz，根據取樣定理，可分析的最高頻率約為取樣頻率的一半：

\[
f_{max} = \frac{F_s}{2} = \frac{22.4k}{2} = 11.2kHz
\]

因此，本系統 FFT 頻譜顯示範圍約為 0 Hz 至 11.2 kHz。

\subsubsection{LED 頻帶分配}

由於 LED 矩陣為 8$\times$8，因此可將頻率分成 8 個頻帶，並分別對應到 LED 矩陣的 8 個欄位。

\begin{table}[H]
\centering
\caption{LED 矩陣頻帶分配}
\begin{tabular}{ccc}
\toprule
LED 欄位 & 頻率範圍 & 顯示意義 \\
\midrule
第 1 欄 & 0--200 Hz & 低頻、鼓聲 \\
第 2 欄 & 200--500 Hz & 低中頻 \\
第 3 欄 & 500 Hz--1 kHz & 人聲低頻 \\
第 4 欄 & 1--2 kHz & 人聲主要區 \\
第 5 欄 & 2--4 kHz & 中高頻 \\
第 6 欄 & 4--6 kHz & 高頻 \\
第 7 欄 & 6--8 kHz & 細節音 \\
第 8 欄 & 8--11.2 kHz & 高頻細節 \\
\bottomrule
\end{tabular}
\end{table}

\subsubsection{LED 高度表示}

FFT 模組會將音訊資料轉換為頻域資料，接著系統將各頻帶能量轉換為 0 到 8 的高度值，再由 LED Matrix 顯示。當低頻能量較強時，左側 LED 欄位高度會增加；當高頻能量較強時，右側 LED 欄位會有較明顯變化。

\begin{table}[H]
\centering
\caption{FFT 能量與 LED 高度對應}
\begin{tabular}{cc}
\toprule
能量大小 & LED 高度 \\
\midrule
很小 & 0--1 格 \\
小 & 2 格 \\
中 & 3--4 格 \\
大 & 5--6 格 \\
很大 & 7--8 格 \\
\bottomrule
\end{tabular}
\end{table}

\subsection{LCD 操作介面設計}

LCD 主要用於顯示目前系統模式與操作提示。顯示內容可設計如下：

\begin{table}[H]
\centering
\caption{LCD 顯示內容規劃}
\begin{tabular}{ccc}
\toprule
狀態 & LCD 第一行 & LCD 第二行 \\
\midrule
主選單 & Audio Recorder & REC PLAY SAVE \\
錄音中 & Recording... & SRAM: xx sec \\
錄音完成 & Record Done & PLAY OK SAVE \\
播放 SRAM & Play SRAM & FFT Display ON \\
儲存 FLASH & Saving FLASH... & Please Wait \\
播放 FLASH & Play FLASH & FFT Display ON \\
錯誤狀態 & Memory Full & Press RESET \\
\bottomrule
\end{tabular}
\end{table}

% \subsubsection{操作規劃}

% \subsubsection{錄音中}

% \subsubsection{錄音完成}

% \subsubsection{播放 SRAM}

% \subsubsection{儲存到 FLASH}

% \subsubsection{播放 FLASH}

% \subsubsection{錯誤顯示}

\subsection{BUTTON / SW 操作規劃}

本系統可使用 BUTTON 與 SW 作為操作介面，控制錄音、播放、儲存與模式選擇。

\begin{table}[H]
\centering
\caption{BUTTON 與 SW 功能規劃}
\begin{tabular}{cc}
\toprule
操作元件 & 功能 \\
\midrule
KEY0 & 返回主選單 / Reset \\
KEY1 & 開始錄音 \\
KEY2 & 播放 / 暫停 \\
KEY3 & 確認 / 存入 FLASH \\
SW0 & 選擇播放來源，0 表示 SRAM，1 表示 FLASH \\
SW1 & 是否開啟 FFT LED 顯示 \\
SW2 & 清除 SRAM 錄音資料 \\
SW3 & 音量或顯示模式切換 \\
\bottomrule
\end{tabular}
\end{table}